\chapter{Math in Bash}

\section{how to run basic math operators}

For you to make basic operations like 10 + 9, you need to use the ``evaluate expression" command.

\lstinputlisting[
    language=bash,
    caption={3-bash.sh},
    captionpos=t,
    label={lst:3-bash},
    breaklines=true]
{../3-bash/3-bash.sh}

\section{Wild cards}

You might think that doing a multiplication is as simple as doing ``expr 10 * 10'' but you'd
be wrong. This is because the asterisk character is a wild card. Wild cards are essentially
special caracters that return specific files depending on their own criteria.
For example, *.log essentially looks for all files that end in .log and returns them to the 
command. this means that a command might perform an operation on all files that end in .log
as oposed of only doing one operation. Using `*' without anything else means to literally use
everything.\\
So how do we solve this issue? we have to use the escape character

\lstinputlisting[
    language=bash,
    caption={3-bash.sh},
    captionpos=t,
    label={lst:3.1-bash},
    breaklines=true]
{../3-bash/3.1-bash.sh}


\section{Math with variables}

Going back to the previous chapter, we can use variables with values and perform operations

\lstinputlisting[
    language=bash,
    caption={3-bash.sh},
    captionpos=t,
    label={lst:3.2-bash},
    breaklines=true]
{../3-bash/3.2-bash.sh}
\pagebreak


\section{Arithmetic expansion}

Now we can continue using expr or use the more modern and prefered way of doing math in general
which is using arithmetic expansions. This don't have space constraints nor require flags.

\lstinputlisting[
    language=bash,
    caption={3-bash.sh},
    captionpos=t,
    label={lst:3.3-bash},
    breaklines=true]
{../3-bash/3.3-bash.sh}

As we'll see later, using the subshell isn't required for if conditionals.